

\section{Bitsets}
Some $O(N^2)$ solutions can be optimized using bitsets. Mainly knapsack problems.
An important point is to use the "popcount" optimization  in target of pragma.

how to initialize:

\begin{lstlisting}
bitset<size> variable_name;
bitset<size> variable_name(DECIMAL_NUMBER);
bitset<size> variable_name("BINARY_STRING")
\end{lstlisting}

Every position can be accessed as a vector, also size must be constant. We get an constant optimization of 32. every bit operation used in integers can be used in bitsets.

Some functions in bitsets:
\begin{center}
\begin{tabular}{ c | l }
  set()  & Set the bit value at all bitset to 1.  \\ 
 reset() & Set the bit value at all bitset to 0.  \\  
 flip() & Flip the bit value at all bitset.\\
 set(idx) & Set the bit value at the given idx to 1.  \\ 
 reset(idx) & Set the bit value at a given idx to 0.  \\  
 flip(idx) & Flip the bit value at the given idx.\\
 count() & Count the quantity of bits set.\\  
 any() & Checks if any bit is set.\\
 all() & Check if all bit is set.\\
 none() & Checks if none bit is set.\\
 size() & Returns the size of the bitset.\\
 to\_string() & Converts bitset to std::string.\\
 to\_ulong() & Converts bitset to unsigned long.\\
 to\_ullong() & Converts bitset to unsigned long long.\\
 \_Find\_first() & return index of first bit set. \\
 \_Find\_next(idx) & return index of first bit set after idx.
\end{tabular}
\end{center}
